% !TEX program = pdflatex
\documentclass[aspectratio=169]{beamer}

\usetheme{Madrid}

\usepackage{amsmath,amssymb,mathtools}
\usepackage{physics}
\usepackage{bm}

\title{Jordan--Wigner Transformation}
\subtitle{From Second Quantization to Qubit Hamiltonians}
\author{Graduate Lecture Notes}
\date{}

\begin{document}

%========================================================
\begin{frame}
\titlepage
\end{frame}

%========================================================
\begin{frame}{Motivation}
\begin{itemize}
\item Quantum computers operate with qubits (Pauli matrices)
\item Many-body physics is naturally expressed in fermionic operators
\item Need mapping:
\[
c^\dagger_p, c_p \;\longrightarrow\; \text{Pauli operators}
\]
\end{itemize}

\begin{block}{Goal}
Map a general second-quantized Hamiltonian to a sum of Pauli strings.
\end{block}
\end{frame}

%========================================================
\begin{frame}{Fermionic Hamiltonian}
General two-body Hamiltonian:
\[
H = \sum_{pq} h_{pq} c^\dagger_p c_q 
+ \frac{1}{2} \sum_{pqrs} V_{pqrs} c^\dagger_p c^\dagger_q c_s c_r
\]

\begin{itemize}
\item One-body term: hopping / kinetic energy
\item Two-body term: interaction
\end{itemize}
\end{frame}

%========================================================
\begin{frame}{Fermionic algebra}
\[
\{c_p, c_q^\dagger\} = \delta_{pq}, \quad
\{c_p, c_q\} = 0
\]

\begin{block}{Challenge}
Qubit operators commute on different sites:
\[
[\sigma_i, \sigma_j] = 0 \quad (i \neq j)
\]

Fermions anticommute!
\end{block}
\end{frame}

%========================================================
\begin{frame}{Jordan--Wigner transformation}
Define mapping:
\[
c_j = \left(\prod_{k<j} Z_k \right) \frac{X_j + i Y_j}{2}
\]

\[
c_j^\dagger = \left(\prod_{k<j} Z_k \right) \frac{X_j - i Y_j}{2}
\]

\begin{block}{Key feature}
Non-local string ensures correct fermionic anticommutation  [oai_citation:0‡Wikipedia](https://en.wikipedia.org/wiki/Jordan%E2%80%93Wigner_transformation?utm_source=chatgpt.com)
\end{block}
\end{frame}

%========================================================
\begin{frame}{Number operator}
\[
n_j = c_j^\dagger c_j = \frac{1}{2}(1 - Z_j)
\]

\begin{itemize}
\item Local operator
\item No Jordan--Wigner string
\end{itemize}
\end{frame}

%========================================================
\begin{frame}{One-body term}
Consider:
\[
c_p^\dagger c_q
\]

Mapping gives:
\[
c_p^\dagger c_q =
\frac{1}{4}
\left(\prod_{k=p}^{q-1} Z_k\right)
(X_p X_q + Y_p Y_q + i X_p Y_q - i Y_p X_q)
\]

\begin{block}{Result}
Becomes sum of Pauli strings with non-local $Z$ string.
\end{block}
\end{frame}

%========================================================
\begin{frame}{Two-body term}
\[
c_p^\dagger c_q^\dagger c_s c_r
\]

\begin{itemize}
\item Each operator contributes:
\begin{itemize}
\item Pauli matrices $(X,Y)$
\item Jordan--Wigner strings
\end{itemize}
\end{itemize}

\medskip

\begin{block}{Result}
Transforms into a sum of Pauli strings of length $\mathcal{O}(N)$.
\end{block}

\medskip

Total number of terms:
\[
\mathcal{O}(N^4)
\]
 [oai_citation:1‡arXiv](https://arxiv.org/pdf/2505.08206?utm_source=chatgpt.com)
\end{frame}

%========================================================
\begin{frame}{Structure after mapping}
General form:
\[
H = \sum_i c_i P_i
\]

where:
\[
P_i = \sigma^{\alpha_1}_1 \otimes \sigma^{\alpha_2}_2 \cdots
\]

\begin{itemize}
\item Pauli strings
\item Non-local due to $Z$ chains
\end{itemize}
\end{frame}

%========================================================
\begin{frame}{Example: Nuclear pairing Hamiltonian}
\[
H = \sum_p \epsilon_p n_p 
- g \sum_{pq} c^\dagger_p c^\dagger_{\bar p} c_{\bar q} c_q
\]

\begin{itemize}
\item Pair creation:
\[
P_p^\dagger = c^\dagger_p c^\dagger_{\bar p}
\]
\end{itemize}
\end{frame}

%========================================================
\begin{frame}{Mapping pairing operators}
Using Jordan--Wigner:

\[
c^\dagger_p c^\dagger_{\bar p}
\;\rightarrow\;
\text{Pauli strings with } X,Y,Z
\]

\medskip

\begin{block}{Important simplification}
Time-reversed pairs often cancel parts of the JW string.
\end{block}

\medskip

Result:
\[
H = \sum_p \epsilon_p \frac{1-Z_p}{2}
- g \sum_{pq} (X_p X_q + Y_p Y_q + \dots)
\]
\end{frame}

%========================================================
\begin{frame}{General two-body interaction}
\[
c^\dagger_p c^\dagger_q c_s c_r
\]

\begin{itemize}
\item Produces:
\begin{itemize}
\item $X X X X$
\item $Y Y X X$
\item mixed terms
\end{itemize}
\end{itemize}

\medskip

\begin{block}{Key observation}
Interaction terms generate entangling Pauli strings.
\end{block}
\end{frame}

%========================================================
\begin{frame}{Scaling and complexity}
\begin{itemize}
\item Number of Pauli terms:
\[
\mathcal{O}(N^4)
\]
\item String length:
\[
\mathcal{O}(N)
\]
\end{itemize}

\medskip

\begin{block}{Implication}
Jordan--Wigner is simple but produces non-local operators.
\end{block}
\end{frame}

%========================================================
\begin{frame}{Physical interpretation}
\begin{itemize}
\item $Z$-strings encode fermionic parity
\item $X,Y$ encode creation/annihilation
\end{itemize}

\medskip

\[
\prod_{k<j} Z_k = (-1)^{\text{# fermions to the left}}
\]

\medskip

\begin{block}{Meaning}
This enforces the Pauli exclusion principle.
\end{block}
\end{frame}

%========================================================
\begin{frame}{Advantages and limitations}
\textbf{Advantages}
\begin{itemize}
\item Simple
\item Exact
\end{itemize}

\textbf{Limitations}
\begin{itemize}
\item Non-local strings
\item Long circuits
\end{itemize}
\end{frame}

%========================================================
\begin{frame}{Summary}
\begin{itemize}
\item Jordan--Wigner maps fermions to qubits
\item General Hamiltonians become sums of Pauli strings
\item Pairing Hamiltonians simplify due to symmetry
\item Two-body interactions produce entangling terms
\end{itemize}

\medskip

\begin{block}{Final insight}
Fermionic statistics $\rightarrow$ non-local qubit operators
\end{block}
\end{frame}

%========================================================
\end{document}


%------------------------------------------------
\section{Explicit Jordan--Wigner Derivation}

\begin{frame}{Goal}
We derive explicitly:
\[
c_p^\dagger c_q^\dagger c_s c_r
\]
in terms of Pauli matrices using the Jordan--Wigner transformation.
\end{frame}

%------------------------------------------------

\begin{frame}{Step 1: JW mapping}
\[
c_j = \left(\prod_{k<j} Z_k \right)\frac{X_j + iY_j}{2}, 
\quad
c_j^\dagger = \left(\prod_{k<j} Z_k \right)\frac{X_j - iY_j}{2}
\]

\medskip

Define:
\[
\sigma_j^+ = \frac{X_j - iY_j}{2}, \quad
\sigma_j^- = \frac{X_j + iY_j}{2}
\]
\end{frame}

%------------------------------------------------

\begin{frame}{Step 2: Write full operator}
\[
c_p^\dagger c_q^\dagger c_s c_r
=
\left(\prod_{k<p} Z_k\right)\sigma_p^+
\left(\prod_{k<q} Z_k\right)\sigma_q^+
\left(\prod_{k<s} Z_k\right)\sigma_s^-
\left(\prod_{k<r} Z_k\right)\sigma_r^-
\]

\medskip

\begin{block}{Key point}
We must carefully combine all $Z$-strings.
\end{block}
\end{frame}

%------------------------------------------------

\begin{frame}{Step 3: Collect parity strings}
Assume ordering:
\[
p < q < r < s
\]

Then:
\[
c_p^\dagger c_q^\dagger c_s c_r
=
\left(\prod_{k<p} Z_k\right)
\left(\prod_{k<q} Z_k\right)
\left(\prod_{k<s} Z_k\right)
\left(\prod_{k<r} Z_k\right)
\cdot
\sigma_p^+ \sigma_q^+ \sigma_s^- \sigma_r^-
\]

\medskip

Combine overlapping strings:
\[
\Rightarrow \prod Z = \prod_{k=p}^{q-1} Z_k \cdot \prod_{k=r}^{s-1} Z_k
\]
\end{frame}

%------------------------------------------------

\begin{frame}{Step 4: Expand ladder operators}
Each operator:
\[
\sigma^\pm = \frac{1}{2}(X \mp iY)
\]

Thus:
\[
\sigma_p^+ \sigma_q^+ \sigma_s^- \sigma_r^-
\]

expands into:
\[
\frac{1}{16} \sum_{\alpha \in \{X,Y\}} (\pm i)^k \;
\sigma_p^\alpha \sigma_q^\beta \sigma_r^\gamma \sigma_s^\delta
\]

\medskip

\begin{block}{Result}
We obtain a sum of $16$ Pauli strings.
\end{block}
\end{frame}

%------------------------------------------------

\begin{frame}{Final expression}
\[
c_p^\dagger c_q^\dagger c_s c_r
=
\frac{1}{16}
\sum_{\alpha,\beta,\gamma,\delta \in \{X,Y\}}
C_{\alpha\beta\gamma\delta}
\left(\prod Z\right)
\sigma_p^\alpha \sigma_q^\beta \sigma_r^\gamma \sigma_s^\delta
\]

\medskip

\begin{itemize}
\item Each term is a Pauli string
\item Coefficients $C_{\alpha\beta\gamma\delta} \in \{\pm 1, \pm i\}$
\end{itemize}

\medskip

\begin{block}{Key result}
Two-body fermionic operators $\Rightarrow$ sum of Pauli strings of length $\mathcal{O}(N)$
\end{block}
\end{frame}


%------------------------------------------------
\section{Bravyi--Kitaev Mapping}

\begin{frame}{Motivation}
Jordan--Wigner:
\begin{itemize}
\item Simple
\item But strings of length $\mathcal{O}(N)$
\end{itemize}

\medskip

Bravyi--Kitaev improves:
\begin{itemize}
\item Locality
\item Scaling
\end{itemize}
\end{frame}

%------------------------------------------------

\begin{frame}{Bravyi--Kitaev idea}
Instead of storing:
\begin{itemize}
\item occupation locally
\end{itemize}

we store:
\begin{itemize}
\item parity + occupation in a binary tree structure
\end{itemize}

\medskip

\begin{block}{Key result}
Operator locality:
\[
\mathcal{O}(\log N)
\]
instead of $\mathcal{O}(N)$
\end{block}
\end{frame}

%------------------------------------------------

\begin{frame}{Comparison}
\begin{center}
\begin{tabular}{lcc}
\toprule
Method & String length & Complexity \\
\midrule
Jordan--Wigner & $\mathcal{O}(N)$ & simple \\
Bravyi--Kitaev & $\mathcal{O}(\log N)$ & more complex \\
\bottomrule
\end{tabular}
\end{center}

\medskip

\begin{block}{Tradeoff}
\begin{itemize}
\item JW: easier to implement
\item BK: better scaling
\end{itemize}
\end{block}
\end{frame}
